\documentclass[12pt,a4paper,oneside]{extreport}
\usepackage{graphicx} % Required for inserting images

% Fontes
\usepackage{fontspec}
\setmainfont{TeX Gyre Pagella}
%\setmainfont{Palatino Linotype}

% Idioma
\usepackage[brazil]{babel}

% Margens
\usepackage[
    left=3cm,
    right=2.5cm,
    top=2.5cm,
    bottom=3cm
]{geometry}

% Espac
\usepackage{setspace}
\onehalfspacing

% Microtipografia
\usepackage{microtype}

%========================================
% Cabeçalhos e Rodapés
%========================================
\usepackage{fancyhdr}

\setlength{\headheight}{16pt}

\pagestyle{fancy}

% Limpa configurações padrão
\fancyhf{}

% Cabeçalho
\fancyhead[R]{{\leftmark}}
\fancyhead[L]{\textbf{Apostila de Programação C}}

% Linha do cabeçalho
\renewcommand{\headrulewidth}{0.4pt}

% Sem linha no rodapé
\renewcommand{\footrulewidth}{0pt}

% Primeira página de cada capítulo
\fancypagestyle{plain}{
    \fancyhf{}
    \fancyfoot[C]{\thepage}
    \renewcommand{\headrulewidth}{0pt}
}

%hyper
\usepackage[hidelinks]{hyperref}


% Pacote codigos
\usepackage{listings}

\lstset{
language=C,
basicstyle=\ttfamily\small,
keywordstyle=\color{blue},
commentstyle=\color{green!60!black},
stringstyle=\color{red},
numbers=left,
numberstyle=\tiny,
frame=single,
breaklines=true
}

% Caixa obs
\usepackage[most]{tcolorbox}

% Cores
\usepackage{xcolor}

\definecolor{titulo}{RGB}{20,55,120}

\setlength{\parskip}{0.4em}


%========================TITULO=======================%

\title{\huge Apostila de Preparação para Programação \textbf{C}}
\author{Rafael Coelho}
%\date{July 2026}

\usepackage{titlesec}

\titleformat{\chapter}
{\Huge\bfseries\color{titulo}}
{\thechapter}
{1em}
{}

%=========================Documento============================%


\begin{document}
\maketitle

%===========================Prefácio=============================%


\chapter*{Prefácio}
\addcontentsline{toc}{chapter}{Prefácio}

A linguagem C ocupa um papel fundamental na formação de profissionais da Computação. Embora linguagens modernas ofereçam mecanismos que abstraem diversos detalhes da implementação, compreender como um programa é executado, como os dados são organizados na memória e como os recursos do computador são utilizados continua sendo uma habilidade indispensável para quem deseja construir uma base sólida em programação.

Esta apostila foi elaborada como material de apoio às monitorias da disciplina de Linguagem C. Seu propósito não é substituir as aulas ou os livros clássicos da área, mas complementar o processo de aprendizagem por meio de explicações objetivas, exemplos comentados, diagramas de memória, exercícios progressivos e observações sobre os erros mais frequentes encontrados pelos estudantes.

Ao longo dos capítulos, a linguagem será apresentada não apenas como uma ferramenta para escrever código, mas como um meio para compreender os fundamentos da programação e o funcionamento interno dos computadores. Sempre que possível, os conceitos serão relacionados a situações práticas e acompanhados de representações visuais que auxiliem na construção de um modelo mental consistente sobre a execução dos programas.

Espera-se que este material sirva como um roteiro de estudos para alunos que desejam consolidar os conhecimentos adquiridos em sala de aula, preparar-se para avaliações e, principalmente, desenvolver uma compreensão mais profunda dos princípios que fundamentam a programação em baixo nível.

\vspace{1cm}

\begin{flushright}
\textit{Bom estudo!}
\end{flushright}

%=====================Como usar=======================%

\chapter*{Como utilizar esta apostila}
\addcontentsline{toc}{chapter}{Como utilizar esta apostila}

%=====================Sumário=======================%

\tableofcontents
\newpage

%===============================================%


\part{Fundamentos de C}

\chapter{Introdução}

Programar é, antes de tudo, resolver problemas. Em sua essência, um programa de computador nada mais é do que uma sequência organizada de instruções capaz de transformar dados de entrada em resultados úteis. Porém, essa definição simples esconde um aspecto fundamental: computadores não compreendem intenções, ideias ou objetivos. Eles executam apenas instruções elementares, em uma ordem rigorosamente definida, realizando milhões ou até bilhões de operações por segundo.

Ao longo da história da Computação, diferentes linguagens de programação foram desenvolvidas para reduzir a distância entre a forma como os seres humanos pensam e a forma como os computadores executam instruções. Linguagens de alto nível, como Python, permitem que problemas complexos sejam descritos por meio de comandos simples e expressivos, abstraindo detalhes relacionados à organização da memória, ao gerenciamento de recursos e à arquitetura do processador. Essa capacidade torna o desenvolvimento de software mais produtivo e acessível, permitindo que o programador concentre seus esforços na lógica da solução.

Entretanto, compreender apenas essas abstrações é conhecer apenas uma parte do processo. Em diversas áreas da Computação, como sistemas operacionais, compiladores, sistemas embarcados, arquitetura de computadores, processamento de alto desempenho e desenvolvimento de drivers, torna-se necessário compreender como os dados são realmente armazenados, manipulados e transferidos pela máquina. É nesse contexto que a linguagem C assume um papel de destaque.

Desenvolvida no início da década de 1970 por Dennis Ritchie, nos Laboratórios Bell, a linguagem C surgiu com o objetivo de oferecer uma alternativa eficiente e portável para o desenvolvimento do sistema operacional UNIX. Desde então, tornou-se uma das linguagens mais influentes da história da Computação, servindo de base para inúmeras outras linguagens modernas e permanecendo amplamente utilizada em aplicações nas quais desempenho, controle de memória e acesso direto ao hardware são requisitos fundamentais.

Estudar C não significa abandonar as facilidades oferecidas pelas linguagens modernas, mas compreender os mecanismos que tornam essas facilidades possíveis. Conceitos como ponteiros, organização da memória, pilha de execução, alocação dinâmica e representação binária dos dados constituem fundamentos presentes, de forma direta ou indireta, em praticamente todas as áreas da Computação. Assim, aprender C representa uma oportunidade de desenvolver uma compreensão mais profunda sobre o funcionamento dos programas e dos computadores.

Ao longo desta apostila, os conceitos serão apresentados de maneira gradual, sempre buscando estabelecer uma relação entre a escrita do código e o comportamento do programa durante sua execução. Sempre que possível, serão utilizados diagramas de memória, exemplos comentados e exercícios progressivos para auxiliar na construção de um modelo mental consistente sobre o funcionamento interno da linguagem.


\section{O que significa programar?}

É um capítulo quase filosófico, mas curto.

Responder perguntas como:

\begin{itemize}
    \item O que é um programa?
    \item O que significa resolver problemas computacionalmente?
    \item O computador "entende" o que escrevemos?
    \item O que diferencia um algoritmo de um programa?
\end{itemize}



\section{Do problema ao algoritmo}

É importante reforçar que programar não é escrever código; escrever código é apenas uma etapa da solução de um problema.

\section{Da linguagem humana à linguagem de máquina}

Essa seção cria a ponte perfeita para C.

Aqui você pode mencionar, sem aprofundar, compiladores, interpretadores e arquiteturas.

Assim, quando chegar o capítulo de compilação, o aluno já terá visto essa ideia.


\section{O surgimento da linguagem C}

Agora sim entra a história.

Não faria uma cronologia extensa.

Explicaria apenas:

Bell Labs;
Dennis Ritchie;
UNIX;
necessidade de uma linguagem eficiente e portátil;
influência sobre linguagens posteriores.

O objetivo é contextualizar, não decorar datas.

\section{Por que estudar C atualmente?}

Essa será uma das seções mais importantes da apostila.

Responder diretamente:

"Por que aprender C se eu já sei Python?"

Explicar que C desenvolve uma compreensão profunda sobre:

memória;
ponteiros;
desempenho;
arquitetura;
sistemas operacionais;
compiladores;
sistemas embarcados.


\section{O primeiro programa em C}

Somente agora aparece o famoso

\begin{lstlisting}
#include <stdio.h>

int main(void)
{
    printf("Olá, mundo!\n");
    return 0;
}
\end{lstlisting}

Mas sem explicar cada linha ainda.

Apenas mostrar que este será o primeiro contato com um programa completo.


\section{Estrutura de um programa em C}

Agora sim analisar linha por linha.

Por exemplo:

'#include'
função main
chaves
instruções
return
comentários

Essa seção já prepara o Capítulo 2.

\section{Compilando o primeiro programa}

Finalizar mostrando o fluxo:

arquivo.c

↓

Pré-processador

↓

Compilador

↓

Assembly

↓

Montador

↓

Arquivo objeto

↓

Linker

↓

Executável


\section*{Resumo}

Exemplo:

Programar é resolver problemas.
Algoritmos descrevem soluções.
C é uma linguagem compilada.
C oferece maior controle sobre os recursos do computador.
O primeiro programa possui uma estrutura básica que será aprofundada nos próximos capítulos.

\section*{Erros comuns}

Exemplo:

\begin{itemize}
    \item esquecer ;
    \item esquecer \{\}
    \item escrever main() incorretamente
    \item esquecer return 0; (explicando que em C moderno ele pode ser omitido em alguns casos, mas mantê-lo é uma boa prática didática).
\end{itemize}

\section*{Exercícios}

Eu dividiria em três níveis.

Fixação

O que é um algoritmo?
Qual a diferença entre linguagem de máquina e linguagem de alto nível?

Aplicação

Explique por que um computador não compreende diretamente um programa escrito em C.

Reflexão

Em quais situações uma linguagem como Python é mais adequada do que C? E quando ocorre o contrário?

%----------------------------------------

\chapter{Variáveis, Tipos e Operadores}

\section{Tipos Primitivos}

\section{Representação em Memória}

\section{Escopo}

%----------------------------------------

\chapter{Controle de Fluxo}

\section{Condicionais}

\section{Laços}

%----------------------------------------

\chapter{Funções}

\section{Declaração}

\section{Protótipos}

\section{Passagem por Valor}


%----------------------------------------

\part{Organização dos Dados}

\chapter{Vetores}

\chapter{Strings}

\chapter{Arrays Multidimensionais}

%========================================

\part{Gerenciamento de Memória}


\chapter{Ponteiros}

% teste codigo
\begin{lstlisting}
int x = 10;
int *p = &x;
\end{lstlisting}


% caixa
\begin{tcolorbox}[title=Importante]

Nunca utilize um ponteiro sem inicializá-lo.

\end{tcolorbox}


\chapter{Organização da Memória}

\section{Segmento de Código}

\section{Stack}

\section{Heap}

\section{Data}

\section{BSS}


\chapter{Alocação Dinâmica}

\section*{Exercícios}


%========================================

\part{Programação Modular}

\chapter{Structs e Enums}

\chapter{Modularização}

\section{Arquivos Header}

\section{Include Guards}

\section{Compilação Separada}

\section{Linkagem}

\chapter{Estruturas de Dados Básicas}



%========================================

\part{Tópicos Avançados}

\chapter{Manipulação Bit a Bit}

\chapter{Recursão}

\chapter{Depuração}

%========================================

\appendix

\part{Apêndices}

\chapter{Perguntas Frequentes}

\chapter{Erros Clássicos em C}

\chapter{Questões Comentadas}

\chapter{Checklist para a Seleção}



\end{document}
